% !TeX root = ../navier-stokes-manuscript.tex
% ============================================================
% 1.7 Singularity and the Derivative Tower
% ============================================================
\subsection{Singularity and the Derivative Tower}\label{sub:SingularityAndTheDerivativeTower}

The vortex sheet broke at the first thing we looked at: velocity had two values across one surface.
A singularity does not have to announce itself that low.
One parcel may keep a finite velocity while the way that velocity is changing loses finite control.
To see that failure, we have to follow the response of the same parcel as far upward as the equation carries it.

Let \(X(a,t)\) be the trajectory of material label \(a\), and write
\[
  \dot X(a,t)=u(X(a,t),t),
  \qquad
  U(a,t):=u(X(a,t),t).
\]
The curve \(U(a,\cdot)\) is the velocity experienced by that parcel.
When this velocity changes, its first change is acceleration.
When the acceleration changes, its next change is jerk.
Each further change is a higher response of the same moving matter.

For this carried sequence, set
\[
  A_r(a,t)
  :=
  \frac{d^r}{dt^r}U(a,t),
  \qquad
  r=0,1,2,\ldots.
\]
On every classical interval,
\[
  A_r(a,t)
  =
  A_r(a,t_0)
  +
  \int_{t_0}^{t}A_{r+1}(a,s)\,ds.
\]
This elementary identity gives the response sequence its finite-time force.
If one rung \(A_r\) escapes every finite bound as \(t\uparrow T<\infty\), the rung \(A_{r+1}\) cannot stay bounded on the remaining interval, because a bounded integrand can change \(A_r\) by only a finite amount.
Applying the same observation again carries the loss of boundedness into every higher material response for as long as those derivatives exist.
Once a rung is genuinely running away at a finite time, the rung above cannot remain still and bounded while producing that change.

The first of these responses is already inside the Navier--Stokes equation.
With
\[
  D_t:=\partial_t+u\cdot\nabla,
\]
the chain rule along the trajectory gives
\[
  A_1(a,t)
  =
  D_tu(X(a,t),t)
  =
  \bigl(-\nabla p+\nu\Delta u\bigr)(X(a,t),t).
\]
Acceleration is the pressure and viscous response of the field occupied by the parcel.
Its jerk is found by differentiating the same equality along the same path:
\[
  A_2(a,t)
  =
  D_t^2u(X(a,t),t)
  =
  D_t\bigl(-\nabla p+\nu\Delta u\bigr)(X(a,t),t).
\]
Higher material differentiation keeps reopening the pressure, viscosity, transport, and spatial variation that generated the response below it.

This carried acceleration must be kept distinct from the time derivative seen at one fixed coordinate.
Write
\[
  V_n:=\partial_t^n u,
  \qquad
  Q_n:=\partial_t^n p.
\]
Then
\[
  A_1
  =
  V_1+(V_0\cdot\nabla)V_0
\]
along the parcel, so \(V_1\) supplies the Eulerian part of material acceleration and transport supplies the other part.
The material jerk contains the next Eulerian time derivative together with the ways the parcel moves through a changing field:
\[
  D_t^2u
  =
  V_2
  +2(V_0\cdot\nabla)V_1
  +(V_1\cdot\nabla)V_0
  +(V_0\cdot\nabla)\bigl((V_0\cdot\nabla)V_0\bigr).
\]
Material response and fixed-coordinate response are two readings of one field, joined by \(D_t=\partial_t+u\cdot\nabla\).

\divider{The Temporal Rungs}

The first four fixed-coordinate time derivatives show how the equation organizes those readings.
The original momentum equation gives
\[
  V_1
  =
  \nu\Delta V_0
  -(V_0\cdot\nabla)V_0
  -\nabla Q_0.
\]
This \(V_1\) is the velocity change at fixed \(x\), and adding transport recovers material acceleration.
Differentiate once in time while holding \(x\) fixed.
The derivative can land on either velocity in the nonlinear product, giving
\[
  V_2
  =
  \nu\Delta V_1
  -(V_1\cdot\nabla)V_0
  -(V_0\cdot\nabla)V_1
  -\nabla Q_1.
\]
Differentiate again, and the middle allocation appears twice:
\[
  V_3
  =
  \nu\Delta V_2
  -(V_2\cdot\nabla)V_0
  -2(V_1\cdot\nabla)V_1
  -(V_0\cdot\nabla)V_2
  -\nabla Q_2.
\]
The fourth rung exposes the next Pascal row:
\[
  \begin{aligned}
  V_4
  ={}&
  \nu\Delta V_3
  -(V_3\cdot\nabla)V_0
  -3(V_2\cdot\nabla)V_1\\
  &-3(V_1\cdot\nabla)V_2
  -(V_0\cdot\nabla)V_3
  -\nabla Q_3.
  \end{aligned}
\]
The coefficients count the ways the time derivatives can divide between the transporting and transported copies of velocity.
At arbitrary order the same product rule reads
\[
  V_{n+1}
  +
  \sum_{k=0}^{n}
  \binom{n}{k}
  (V_k\cdot\nabla)V_{n-k}
  +
  \nabla Q_n
  =
  \nu\Delta V_n,
  \qquad
  \nabla\cdot V_n=0.
\]
Each \(V_{n+1}\) is tied to lower time rows through transport, to spatial curvature through viscosity, and to pressure at the same time order.
The recurrence is an identity on the classical solution, and every row derives from that shared evolution.

\divider{Pressure and the Spatial Rungs}

Pressure and spatial derivatives have entered every temporal row from the beginning.
Taking the divergence of the original momentum equation gives
\[
  -\Delta p
  =
  \partial_i\partial_j(u_i u_j),
\]
where repeated spatial indices are summed.
After \(n\) time derivatives,
\[
  -\Delta Q_n
  =
  \partial_i\partial_j
  \sum_{k=0}^{n}
  \binom{n}{k}
  V_{k,i}V_{n-k,j}.
\]
The whole-space decay normalization fixes the spatial constant left open by the elliptic equation.
Each pressure row is generated by the velocity rows across the field, so a pressure response at one parcel can depend on velocity far from that parcel.

Spatial differentiation opens the momentum equation in a second direction.
One derivative in \(x_\ell\) gives
\[
  \begin{aligned}
  \partial_t(\partial_\ell u)
  &+
  (u\cdot\nabla)(\partial_\ell u)
  +
  ((\partial_\ell u)\cdot\nabla)u
  +
  \nabla(\partial_\ell p)\\
  &=
  \nu\Delta(\partial_\ell u).
  \end{aligned}
\]
The added transport term records the fact that the transporting velocity also varies across neighbouring positions.
A second spatial derivative gives both cross allocations:
\[
  \begin{aligned}
  \partial_t(\partial_{\ell m}u)
  &+
  (u\cdot\nabla)(\partial_{\ell m}u)
  +
  ((\partial_\ell u)\cdot\nabla)(\partial_m u)\\
  &+
  ((\partial_m u)\cdot\nabla)(\partial_\ell u)
  +
  ((\partial_{\ell m}u)\cdot\nabla)u
  +
  \nabla(\partial_{\ell m}p)\\
  &=
  \nu\Delta(\partial_{\ell m}u).
  \end{aligned}
\]
Every further temporal or spatial differentiation repeats this nesting and creates the corresponding mixed response row.

Using a multi-index \(\alpha\), write
\[
  J_{n,\alpha}(x,t)
  :=
  \partial_t^n\partial_x^\alpha u(x,t),
\]
and collect the velocity readings at that spacetime point as
\[
  \mathcal J^\infty u(x,t)
  :=
  \bigl(
    J_{n,\alpha}(x,t)
  \bigr)_{n\geq0,\,\alpha\in\mathbb N_0^3}.
\]
This is the derivative tower.
Its rungs are simultaneous readings of one velocity field, its material acceleration and jerk are combinations of those readings along a trajectory, and its pressure rows remain attached through the Poisson equation and incompressibility.

The tower can contain lawful zeroes.
The zero solution has every rung equal to zero, a smooth steady flow has \(V_n=0\) for all \(n\geq1\), and a changing response can pass through zero at an instant.
The equation decides whether a zero or nonzero value fits the surrounding field.
The finite-time escalation above begins only when a particular material rung actually escapes every finite bound.

A finite terminal velocity leaves all of these higher readings unsettled.
Smooth arrival requires each finite temporal, spatial, and mixed rung to extend with one finite value compatible with the equations.
The board-like jump is the lowest visible failure of that requirement; the opened tower shows where less visible failures can occur.
